% Extracción quirúrgica del propietario V2 05c: líneas 99--335.
% El resto permanece en este archivo como procedencia, pero no se publica.
\iffalse
\chapter[Macrostructure de jauge interne]{Macrostructure de jauge interne engendrée par le continu discret}
\label{chap:gauint-macroestructura}

\begin{thesisbox}
La macrostructure de ce chapitre se construit uniquement à partir d'une
restriction survivante de \(\mathcal C_{\rm cont}^{\rm disc}\). Les nombres
\(4\), \(6\) et \(8\), l'écran, les quatre involutions, la mesure
projective, l'opérateur de transfert, le mode d'incidence et le terminal
cellulaire se déduisent au sein de cette restriction. Aucune évaluation ultérieure
ne fait partie du domaine ni des flèches qui suivent.
\end{thesisbox}

\section{Le germe projectif et la dérivation interne de \texorpdfstring{\(4/6/8\)}{4/6/8}}

Soit \(V_{\rm TPK}=\mathbf F_3^2\), avec la base ordonnée héritée des deux
feuillets de l'état TPK. Nous définissons l'ensemble des directions
\begin{equation}\label{eq:gauint-cuatro-direcciones}
 I_{\rm gau}:=\mathbf P^1(\mathbf F_3)
 =\{L_\infty,L_0,L_1,L_{-1}\},
\end{equation}
où
\[
 L_\infty=[1:0],\qquad L_0=[0:1],\qquad
 L_1=[1:1],\qquad L_{-1}=[1:-1].
\]
Par conséquent,
\begin{equation}\label{eq:gauint-cardinal-cuatro}
 |I_{\rm gau}|=\frac{3^2-1}{3-1}=4.
\end{equation}
On n'introduit pas un ensemble de quatre indices : il apparaît comme le quotient des
huit vecteurs non nuls par l'action de \(\mathbf F_3^\times\).

Les incidences non orientées sont
\begin{equation}\label{eq:gauint-seis-incidencias}
 E_{\rm gau}:=\binom{I_{\rm gau}}2,
 \qquad |E_{\rm gau}|=\binom42=6,
 \qquad Q_{\rm gau}:=\mathbf F_3^{E_{\rm gau}}.
\end{equation}
L'écran orienté double chaque direction, et non chaque incidence :
\begin{equation}\label{eq:gauint-ocho-caras}
 \mathscr B_{\rm gau}:=I_{\rm gau}\times\{+,-\},
 \qquad |\mathscr B_{\rm gau}|=2|I_{\rm gau}|=8.
\end{equation}
Ainsi, \(4\), \(6\) et \(8\) proviennent respectivement de directions, de couples de
directions et de faces orientées d'un même germe.

Pour chaque \(L\in I_{\rm gau}\), soit \(\vartheta_L\) l'involution de
l'écran qui échange \((L,+)\) et \((L,-)\) et laisse fixes les six autres
faces. Alors
\begin{equation}\label{eq:gauint-cuatro-involuciones}
 \vartheta_L^2=I,
 \qquad
 \vartheta_L\vartheta_{L'}=\vartheta_{L'}\vartheta_L
 \quad(L\ne L').
\end{equation}
Les quatre involutions ont été obtenues avant de définir une mesure ou un opérateur.

\section{Incidence, mode interne et témoin non constant}

L'application d'incidence est
\begin{equation}\label{eq:gauint-incidencia-b}
 B:\mathbf F_3^{I_{\rm gau}}\longrightarrow Q_{\rm gau},
 \qquad
 (Ba)_{\{L,L'\}}:=a_L+a_{L'}.
\end{equation}
Si
\[
 s(q):=\sum_{e\in E_{\rm gau}}q_e\in\mathbf F_3,
\]
chaque direction apparaît exactement trois fois dans la somme des six
incidences, et donc
\begin{equation}\label{eq:gauint-incidencia-invariante}
 s(Ba)=3\sum_{L\in I_{\rm gau}}a_L=0,
 \qquad
 s(q+Ba)=s(q).
\end{equation}

Sur \(Q_{\rm gau}\), on utilise exclusivement la moyenne de comptage
normalisée. Le mode d'incidence est la fonction rationnelle
\begin{equation}\label{eq:gauint-wc-definicion}
 W_c(q):=\mathbf 1_{\{s(q)=0\}}-\frac13.
\end{equation}
Comme \(s:Q_{\rm gau}\to\mathbf F_3\) est surjective, chaque fibre a
\(3^5\) éléments. On en obtient, par comptage exact,
\begin{equation}\label{eq:gauint-wc-momentos}
 \mathbb E_{Q}W_c=0,
 \qquad
 \mathbb E_{Q}W_c^2=\frac29,
 \qquad
 \mathbb E_{Q}W_c^3=\frac{2}{27}.
\end{equation}
En outre, \eqref{eq:gauint-incidencia-invariante} donne
\begin{equation}\label{eq:gauint-wc-invariancia}
 W_c(q+Ba)=W_c(q).
\end{equation}
Donc \(W_c\) est non nul, centré et invariant par l'incidence
engendrée en interne.

\fi
\section{Incidence, courbure de retenue et domaine commun de naturalité}

Pour \(x\in\mathcal C_{\rm cont}^{\rm disc}\), notons
\(Y_m(x)\) l'ensemble fini des préfixes de jauge de profondeur \(m\) qui
conservent simultanément
\begin{equation*}
\begin{gathered}
 \text{fibre, relations d’extension, nombres, orientation, retenue, mémoire,}\\
 \text{frontière, route, multisection, survie, terminal et lecteur}.
\end{gathered}
\end{equation*}
Il existe une application de collier
\begin{equation}\label{eq:gauint-collar}
 q_m:Y_m(x)\longrightarrow Q_{\rm gau}
\end{equation}
qui enregistre les six incidences sans effacer l'histoire qui les produit.
Pour \(\ell\ge m\), la troncature s'écrit
\[
 \tau_{\ell,m}:Y_\ell(x)\longrightarrow Y_m(x).
\]

Pour \(y,z\in Y_m(x)\), soit
\begin{equation}\label{eq:gauint-todas-extensiones}
 \operatorname{Ext}^{\rm adm}_{m}(y,z)
\end{equation}
l'ensemble de \emph{toutes} les extensions TPK admissibles de \(y\) à \(z\).
Une extension appartient à cet ensemble seulement si elle transporte la fibre entière,
les relations \(\operatorname{Ext}_\Gamma\), la retenue native et opératorielle,
l'orientation, la mémoire accumulée, la frontière, le chemin et la condition de
survie. Nous définissons
\begin{equation}\label{eq:gauint-conteo-kernel}
 a_m(y,z):=\#\operatorname{Ext}^{\rm adm}_m(y,z),
 \qquad
 d_m(y):=\sum_{z\in Y_m(x)}a_m(y,z),
\end{equation}
et, lorsque \(d_m(y)>0\),
\begin{equation}\label{eq:gauint-kernel-normalizado}
 k_m(y,z):=\frac{a_m(y,z)}{d_m(y)}.
\end{equation}
Ainsi, \(k_m\) est un comptage normalisé d'extensions complètes, non une
fonction introduite après la projection de l'état.

La restriction \(\mathcal C_{\rm cont}^{\rm gau}\) est formée des
histoires survivantes \(x\) pour lesquelles les conditions
suivantes sont satisfaites à toute profondeur.
\begin{enumerate}
 \item Les ensembles \(Y_m(x)\) sont non vides, \(d_m(y)>0\), et les
 troncatures \(\tau_{\ell,m}\) sont surjectives.
 \item Il existe une multisection non vide de poids rationnels strictement
 positifs \(\mu_m:Y_m(x)\to\mathbf Q_{>0}\) telle que
 \begin{equation}\label{eq:gauint-medida-proyectiva}
  \sum_{y\in Y_m}\mu_m(y)=1,
  \qquad
  \mu_m(y)=\sum_{\tau_{\ell,m}\widetilde y=y}\mu_\ell(\widetilde y).
 \end{equation}
 \item Chaque \(\mu_m\) est stationnaire :
 \begin{equation}\label{eq:gauint-estacionariedad}
  \sum_{y\in Y_m}\mu_m(y)k_m(y,z)=\mu_m(z).
 \end{equation}
 \item La marginale de \((q_m)_*\mu_m\) est le comptage uniforme de
 \(Q_{\rm gau}\). Cette condition rend littéralement applicables les calculs
 de \eqref{eq:pdfv2-cor-wc-momentos} au relèvement \(W_c\circ q_m\).
 \item La lumpabilité directe est satisfaite
 \begin{equation}\label{eq:gauint-lump-directa}
  \sum_{\tau_{\ell,m}\widetilde z=z}
   k_\ell(\widetilde y,\widetilde z)
  =k_m(\tau_{\ell,m}\widetilde y,z)
 \end{equation}
 pour tout \(\widetilde y\in Y_\ell\) et \(z\in Y_m\).
 \item Pour le noyau adjoint
 \begin{equation}\label{eq:gauint-kernel-adjunto}
  k_m^*(z,y):=\frac{\mu_m(y)k_m(y,z)}{\mu_m(z)},
 \end{equation}
 la lumpabilité adjointe est satisfaite
 \begin{equation}\label{eq:gauint-lump-adjunta}
  \sum_{\tau_{\ell,m}\widetilde y=y}
   k_\ell^*(\widetilde z,\widetilde y)
  =k_m^*(\tau_{\ell,m}\widetilde z,y).
 \end{equation}
 \item Les catégories de préfixes possèdent la racine TPK comme objet initial,
 de sorte que leur nerf possède la contraction terminale explicite de
 \eqref{eq:gauint-sdr} ci-dessous.
\end{enumerate}
Ces conditions définissent le domaine typé ; en dehors de celui-ci, on ne publie pas une
mesure ou un opérateur naturel par simple déclaration.

\section{Mesure projective, noyau et naturalité directe et adjointe}

Une branche de la multisection \(\{\mu_m\}_m\) étant fixée, soit
\begin{equation}\label{eq:gauint-hm}
 H_m:=\operatorname{Fun}(Y_m,\mathbf Q),
 \qquad
 \langle f,g\rangle_m:=\sum_{y\in Y_m}\mu_m(y)f(y)g(y).
\end{equation}
L'opérateur engendré par toutes les extensions est
\begin{equation}\label{eq:gauint-km}
 (K_mf)(y):=\sum_{z\in Y_m}k_m(y,z)f(z).
\end{equation}
Les équations de Markov et de stationnarité donnent
\begin{equation}\label{eq:gauint-markov-estacionario}
 K_m\mathbf1=\mathbf1,
 \qquad
 K_m^*\mathbf1=\mathbf1,
 \qquad
 \|K_mf\|_m\le\|f\|_m.
\end{equation}

Le relèvement et l'espérance de troncature sont
\begin{equation}\label{eq:gauint-lift-esperanza}
 J_{\ell,m}f:=f\circ\tau_{\ell,m},
 \qquad
 E_{m,\ell}:=J_{\ell,m}^*.
\end{equation}
La projectivité de \(\mu\) implique
\begin{equation}\label{eq:gauint-ej-identidad}
 E_{m,\ell}J_{\ell,m}=I_{H_m}.
\end{equation}
Les deux conditions de lumpabilité établissent, respectivement,
\begin{equation}\label{eq:gauint-naturalidad-directa-adjunta}
 K_\ell J_{\ell,m}=J_{\ell,m}K_m,
 \qquad
 K_\ell^*J_{\ell,m}=J_{\ell,m}K_m^*.
\end{equation}
La seconde égalité ne se déduit pas de la première sans la condition adjointe :
les deux sont enregistrées dans le domaine.

On définit l'opérateur interne positif
\begin{equation}\label{eq:gauint-transferencia-t}
 T_m:=K_m^*K_m.
\end{equation}
Alors
\begin{equation}\label{eq:gauint-t-positivo-natural}
 \langle f,T_mf\rangle_m=\|K_mf\|_m^2\ge0,
 \qquad
 T_\ell J_{\ell,m}=J_{\ell,m}T_m.
\end{equation}

\section{Loi commune de \(8\) faces et \(4\) carrés de Fubini}

Sur \(Y_m^{\mathscr B_{\rm gau}}\), on définit la masse rationnelle
\begin{equation}\label{eq:gauint-ley-ocho-caras}
 \Omega_m^{(8)}(\mathbf y)
 :=\sum_{z\in Y_m}\mu_m(z)
 \prod_{(L,\varepsilon)\in\mathscr B_{\rm gau}}
 k_m(z,y_{L,\varepsilon}).
\end{equation}
C'est une probabilité parce que chaque facteur est markovien. Toutes les faces proviennent
du même état latent \(z\) et du même noyau ; ce ne sont pas huit mesures
indépendantes ajoutées.

Pour des fonctions \(f_{L,+},f_{L,-}\in H_m\), Fubini fini donne
\begin{equation}\label{eq:gauint-fubini-ocho}
\begin{aligned}
 &\sum_{\mathbf y}\Omega_m^{(8)}(\mathbf y)
  \prod_{L\in I_{\rm gau}}
  f_{L,+}(y_{L,+})f_{L,-}(y_{L,-})\\
 &\qquad=
 \sum_{z\in Y_m}\mu_m(z)
 \prod_{L\in I_{\rm gau}}
 (K_mf_{L,+})(z)(K_mf_{L,-})(z).
\end{aligned}
\end{equation}
En prenant \(f_{L,+}=f_{L,-}=f_L\), les quatre
carrés apparaissent simultanément
\begin{equation}\label{eq:gauint-cuatro-cuadrados}
 \sum_{z\in Y_m}\mu_m(z)
 \prod_{L\in I_{\rm gau}}(K_mf_L(z))^2.
\end{equation}
La marginale de tout couple opposé satisfait
\begin{equation}\label{eq:gauint-marginal-t}
 \sum_{y_+,y_-}\Omega_{m,L}^{(2)}(y_+,y_-)
 f(y_+)g(y_-)
 =\langle K_mf,K_mg\rangle_m
 =\langle f,T_mg\rangle_m.
\end{equation}
La commutativité du produit dans \eqref{eq:gauint-ley-ocho-caras} prouve
\begin{equation}\label{eq:gauint-involuciones-ley}
 (\vartheta_L)_*\Omega_m^{(8)}=\Omega_m^{(8)}
 \qquad(L\in I_{\rm gau}).
\end{equation}

\section{Opérateur interne exact et témoin propre}

La projection sur le secteur constant et son complément sont
\begin{equation}\label{eq:gauint-pq}
 P_{\Omega,m}f:=\langle\mathbf1,f\rangle_m\mathbf1,
 \qquad
 Q_m:=I-P_{\Omega,m}.
\end{equation}
L'opérateur cellulaire interne est défini par
\begin{equation}\label{eq:gauint-d-hmt}
 D_m^{\rm HMT}:=3Q_m.
\end{equation}
Le coefficient \(3\) est le cardinal du corps des résidus utilisé dans
\eqref{eq:pdfv2-cor-proyectores-cuatro-direcciones} ; il ne représente pas une
échelle ajoutée.
Pour le relèvement
\[
 W_{c,m}:=W_c\circ q_m
\]
on a, par \eqref{eq:pdfv2-cor-wc-momentos},
\begin{equation}\label{eq:gauint-w-eigen}
 P_{\Omega,m}W_{c,m}=0,
 \qquad
 D_m^{\rm HMT}W_{c,m}=3W_{c,m},
 \qquad
 \|W_{c,m}\|_m^2=\frac29.
\end{equation}
Ainsi, le secteur complémentaire possède un témoin non nul construit par
l'incidence des six arêtes.

\section{Terminal cellulaire et rétracte fort}

Soit \(C_*(N\mathsf Y_m;A_n)\) le complexe normalisé du nerf de la
catégorie des préfixes, et soit \(y_*\) sa racine initiale. L'inclusion et
l'augmentation sont
\[
 \iota_m:A_n[0]\longrightarrow C_*(N\mathsf Y_m;A_n),
 \quad \iota_m(1)=[y_*],
 \qquad
 \epsilon_m:C_*(N\mathsf Y_m;A_n)\longrightarrow A_n[0].
\]
L'objet initial fournit la dégénérescence supplémentaire
\begin{equation}\label{eq:gauint-homotopia-celular}
 H_m[y_0\to\cdots\to y_q]
 :=[y_*\to y_0\to\cdots\to y_q],
\end{equation}
avec la convention normalisée lorsqu'une dégénérescence apparaît. On vérifie
\begin{equation}\label{eq:gauint-sdr}
 \epsilon_m\iota_m=I,
 \qquad
 \partial H_m+H_m\partial=I-\iota_m\epsilon_m,
 \qquad
 H_m^2=H_m\iota_m=\epsilon_mH_m=0.
\end{equation}
Ce terminal cellulaire est distinct de la projection constante
\(P_{\Omega,m}\) : le premier contracte le nerf et la seconde sépare les
secteurs de l'espace des fonctions. Les deux sont conservés.

\iffalse
\section{Les treize coordonnées de la restriction de jauge}

Pour chaque \((m,n)\), on définit
\begin{equation}\label{eq:gauint-d-trece}
 D_{{\rm gau},m,n}(x):=
 (\mathcal F,\operatorname{Ext}_\Gamma,\operatorname{Rel},
 \mathcal N,\operatorname{or},\operatorname{Carr},\operatorname{Mem},
 \partial,\gamma,\operatorname{Msect},\operatorname{Surv},
 \operatorname{Term},\mathcal L)_{{\rm gau},m,n}(x),
\end{equation}
avec le contenu suivant.
\begin{enumerate}
 \item \(\mathcal F=(V_{\rm TPK},I_{\rm gau},E_{\rm gau},
 Q_{\rm gau},\mathscr B_{\rm gau},Y_m,H_m)\).
 \item \(\operatorname{Ext}_\Gamma\) contient tous les ensembles
 \(\operatorname{Ext}^{\rm adm}_m(y,z)\), les troncatures, les relèvements,
 les espérances et leurs compositions.
 \item \(\operatorname{Rel}\) contient \(B\), \(sB=0\), les quatre
 involutions, Markov, la stationnarité et les deux lumpabilités.
 \item \(\mathcal N=(3,4,6,8,m,n,|Y_m|,a_m,d_m)\) ; ces nombres ne
 remplacent aucune des autres coordonnées.
 \item \(\operatorname{or}\) conserve les faces \(+/-\), les quatre
 inversions \(\vartheta_L\) et le renversement de chaque chemin.
 \item \(\operatorname{Carr}\) contient la retenue complète de chaque
 extension comptée dans \(a_m(y,z)\), y compris la retenue déjà propagée.
 \item \(\operatorname{Mem}\) conserve le préfixe complet, l'état
 commun \(z\) de \(\Omega_m^{(8)}\) et la mémoire de ses huit extensions.
 \item \(\partial=(q_m,\partial_{N\mathsf Y_m},Q_m,H_m)\) conserve
 la frontière de collier, la frontière cellulaire et le complément.
 \item \(\gamma\) est le simplexe orienté du nerf qui enregistre le chemin
 et non seulement ses extrémités.
 \item \(\operatorname{Msect}\) est la famille de toutes les mesures
 stationnaires projectives compatibles et de toutes les extensions qui
 interviennent dans les comptages ; aucune branche rétrospective n'est sélectionnée.
 \item \(\operatorname{Surv}\) enregistre la non-vacuité, la profondeur
 arbitraire, la projectivité, les deux lumpabilités et la persistance de la
 marginale uniforme de collier.
 \item \(\operatorname{Term}=(\iota_m,\epsilon_m,H_m;
 P_{\Omega,m},Q_m)\) conserve séparément les deux terminaux.
 \item \(\mathcal L=(\Omega_m^{(8)},K_m,T_m,D_m^{\rm HMT},W_{c,m})\)
 est le lecteur interne postérieur à la construction de l'histoire.
\end{enumerate}

Pour \(\ell\ge m\) et \(n'\ge n\), toutes les coordonnées satisfont
\begin{equation}\label{eq:gauint-naturalidad-trece}
 u^{\rm gau}_{(\ell,n'),(m,n)}D_{{\rm gau},\ell,n'}(x')
 =D_{{\rm gau},m,n}(\tau_{\ell,m}x').
\end{equation}

\section{Assembleur, publicateur et chaîne non ablative}

Soit
\begin{equation}\label{eq:gauint-graph-d}
 \mathcal G_{\rm gau}:=\operatorname{Graph}(D_{\rm gau}),
 \qquad
 \operatorname{Res}_{\rm gau}(x):=(x,D_{\rm gau}x).
\end{equation}
La multisection dépendante \(\chi_{\rm gau}(x)\) contient toutes les branches
compatibles de collier, de mesure et d'extension. Nous définissons
\begin{equation}\label{eq:gauint-x-dependiente}
 X_{\chi,{\rm gau}}
 :=\{(\chi_{\rm gau}(x),x,D_{\rm gau}x):
 x\in\mathcal C_{\rm cont}^{\rm gau}\},
\end{equation}
\begin{equation}\label{eq:gauint-prx}
 \operatorname{pr}_{X,{\rm gau}}(x,D_{\rm gau}x)
 :=(\chi_{\rm gau}(x),x,D_{\rm gau}x).
\end{equation}

L'assembleur brut
\begin{equation}\label{eq:gauint-ensamblador-crudo}
 a_{\rm gau}:X_{\chi,{\rm gau}}\longrightarrow
 \mathscr M_{\rm gau}^{\rm amb}
\end{equation}
produit la famille complète
\[
 \bigl(I,E,\mathscr B,B,W_c;
 Y_m,\mu_m,k_m,J_{\ell,m},E_{m,\ell};
 K_m,T_m,\Omega_m^{(8)},P_{\Omega,m},Q_m,D_m^{\rm HMT};
 \iota_m,\epsilon_m,H_m\bigr)_{m,n}.
\]
Le macro-objet conserve son entrée par
\begin{equation}\label{eq:gauint-graph-a}
 \mathfrak M_{\rm gau}:=\operatorname{Graph}(a_{\rm gau}),
 \qquad
 \mathcal A_{\rm gau}(z):=(z,a_{\rm gau}z).
\end{equation}

Le publicateur brut
\begin{equation}\label{eq:gauint-publicador-crudo}
 p_{\rm gau}:\mathfrak M_{\rm gau}\longrightarrow
 \mathscr Y_{\rm gau}
\end{equation}
publie conjointement les identités
\[
 \begin{gathered}
 |I|=4,quad |E|=6,quad |\mathscr B|=8,quad sB=0,
 \quad W_c(q+Ba)=W_c(q),\\
 K_\ell J=JK_m,quad K_\ell^*J=JK_m^*,
 \quad T_m=K_m^*K_m,\\
 (\vartheta_L)_*\Omega_m^{(8)}=\Omega_m^{(8)},
 \quad \text{les quatre factorisations quadratiques de Fubini},\\
 D_m^{\rm HMT}=3Q_m,quad D_m^{\rm HMT}W_{c,m}=3W_{c,m},
 \quad \partial H_m+H_m\partial=I-\iota_m\epsilon_m.
 \end{gathered}
\]
Enfin,
\begin{equation}\label{eq:gauint-graph-p}
 \mathcal C_{\rm gau}^{\rm HMT}:=\operatorname{Graph}(p_{\rm gau}),
 \qquad
 \mathcal P_{\rm gau}(M):=(M,p_{\rm gau}M).
\end{equation}
La chaîne complète est
\begin{equation}\label{eq:gauint-cadena-completa}
 \mathcal C_{\rm cont}^{\rm disc}\supset
 \mathcal C_{\rm cont}^{\rm gau}
 \xrightarrow{\operatorname{Res}_{\rm gau}}\mathcal G_{\rm gau}
 \xrightarrow{\operatorname{pr}_{X,{\rm gau}}}X_{\chi,{\rm gau}}
 \xrightarrow{\mathcal A_{\rm gau}}\mathfrak M_{\rm gau}
 \xrightarrow{\mathcal P_{\rm gau}}\mathcal C_{\rm gau}^{\rm HMT}.
\end{equation}
Chaque flèche possède un inverse sur son image donné par une projection de
coordonnées. Aucune sortie ne remplace son histoire génératrice.

\section{Théorème de réalisation de jauge interne}

\begin{theorem}[Réalisation interne complète et naturelle]
\label{thm:gauint-realizacion}
Pour tout \(x\in\mathcal C_{\rm cont}^{\rm gau}\), la chaîne
\eqref{eq:gauint-cadena-completa} construit, à partir de la structure
discrète du continu, une macrostructure naturelle avec quatre directions,
six incidences, huit faces, quatre involutions, une mesure projective, un
noyau de comptage complet, son transfert positif, une loi commune de huit
faces, un témoin d'incidence non constant et un terminal cellulaire. Les
treize coordonnées de \eqref{eq:gauint-d-trece} restent récupérables en
sortie.
\end{theorem}

\begin{proof}
Les équations \eqref{eq:gauint-cuatro-direcciones}--
\eqref{eq:gauint-ocho-caras} déduisent \(4/6/8\). La somme des incidences
\eqref{eq:gauint-incidencia-invariante} prouve l'invariance de \(W_c\), et
le comptage uniforme prouve \eqref{eq:gauint-wc-momentos}. Les extensions
complètes produisent \(k_m\) par \eqref{eq:gauint-kernel-normalizado} ; les
conditions explicites du domaine prouvent Markov, la stationnarité et les deux
égalités de naturalité. Donc \(T_m=K_m^*K_m\) est positif et naturel.
Fubini appliqué une seule fois à \eqref{eq:gauint-ley-ocho-caras} donne les quatre
formes quadratiques sans changer de mesure. La marginale uniforme centre
\(W_{c,m}\), d'où \eqref{eq:gauint-w-eigen}. L'objet initial du nerf
donne la dégénérescence supplémentaire et l'identité SDR \eqref{eq:gauint-sdr}. Enfin,
les trois constructions par graphe retiennent leurs entrées comme premières
coordonnées, donc la composition n'oublie aucune des treize entrées.
\end{proof}

\section{Provenance et falsificateurs}

\paragraph{Provenance.}
Le germe TPK, l'écran orienté, les extensions survivantes, la
mesure commune et la contraction terminale ont le statut
\texttt{ARCHITECTURE\_AUTORALE\_PRÉEXISTANTE}. Le transfert par comptage,
le mode d'incidence et le secteur exact \(3Q\) sont présentés comme
\texttt{RÉSULTAT\_RÉCUPÉRÉ}. La dérivation explicite
\(\mathbf P^1(\mathbf F_3)\to4/6/8\), les deux conditions de lumpabilité et
la mise en paquet non ablative par graphes ont le statut
\texttt{FORMALISATION\_NOUVELLE} ; elles formalisent l'architecture antérieure sans
revendiquer une origine conceptuelle nouvelle.

\begin{ablation}[Falsificateurs internes de la branche de jauge]
\label{abl:gauint-falsadores}
La construction est falsifiée au niveau indiqué si l'un quelconque des
faits suivants se produit.
\begin{enumerate}
 \item On remplace \(I_{\rm gau}\), \(E_{\rm gau}\) ou
 \(\mathscr B_{\rm gau}\) par une liste sans les applications qui les déduisent.
 \item \(sB=0\) échoue ou \(W_c\) cesse d'être invariant et non constant.
 \item Le comptage \(a_m(y,z)\) omet la retenue, l'orientation, la mémoire, la frontière ou
 le chemin : dans ce cas, on a compté un autre objet.
 \item La projectivité, la stationnarité ou l'une des deux
 lumpabilités échoue ; on ne peut alors publier la naturalité correspondante.
 \item Les huit faces ne proviennent pas d'une unique \(\mu_m\) et d'un unique
 \(k_m\), ou l'une des \(\vartheta_L\) ne préserve pas \(\Omega_m^{(8)}\).
 \item \(P_{\Omega,m}W_{c,m}\ne0\) o
 \(D_m^{\rm HMT}W_{c,m}\ne3W_{c,m}\).
 \item L'identité SDR \(\partial H+H\partial=I-\iota\epsilon\) échoue.
 \item L'une des treize coordonnées est omise : \emph{objet substitué ;
 conclusion non transférable}.
\end{enumerate}
Ces falsificateurs contrôlent l'appartenance et la naturalité ; ils ne dégradent pas une
histoire qui satisfait toutes les identités.
\end{ablation}
\fi
